\documentclass[tikz,border=4pt]{standalone}
\usepackage{ctex}
\usepackage{amsmath}
\usetikzlibrary{calc, arrows.meta}

\begin{document}
\begin{tikzpicture}[
  scale=1.0, thick,
  vec/.style={-{Stealth[length=7pt, width=4pt]}, thick},
]

% 以公共起点 O 画四条向量，纵向排列便于对比

% 基向量 a（长度单位：2cm，方向向右上）
\coordinate (O) at (0, 0);
\def\ax{2.0} \def\ay{1.0}

% vec a：蓝色，实线
\draw[vec, blue]   (O) -- ++ (\ax,   \ay)   node[right] {$\vec{a}$};

% vec 2a：黑色，实线（长度 ×2）
\draw[vec, black]  (O) -- ++ ({2*\ax}, {2*\ay}) node[right] {$2\vec{a}$};

% vec -a：红色（反向）
\draw[vec, red]    (O) -- ++ ({-\ax},  {-\ay})  node[left]  {$-\vec{a}$};

% vec 1/2 a：green!60!black，长度 ×1/2
\draw[vec, green!60!black] (O) -- ++ ({0.5*\ax}, {0.5*\ay})
  node[below right, xshift=1pt] {$\dfrac{1}{2}\vec{a}$};

% 原点标记
\filldraw[black] (O) circle (1.5pt) node[below left] {$O$};

% 说明文字
\node[align=left, font=\small] at (3.5, -1.2) {
  同起点，同/反方向，\\
  长度按 $|\lambda|$ 伸缩
};

\end{tikzpicture}
\end{document}
